\section{Discussion}
\label{sec:discussion}

\subsection{Il grado di gossip come compromesso esplicito}
A $N$ fissato, un $k$ più alto porta più aggiornamenti per round e, nella
maggior parte delle configurazioni testate ($N=5$ e $N=8$, sia in locale sia
su AWS), un'accuratezza global-test più alta a parità di $H$ — coerente con
l'attesa che un mixing più denso tra i pesi dei worker riduca la divergenza
verso soluzioni funzionalmente diverse. A $N=3$ l'effetto è invece piatto o
leggermente negativo in entrambi gli ambienti: con solo due vicini possibili,
passare da $k=1$ a $k=2$ non introduce una differenza qualitativa nella
topologia di comunicazione, e il rumore tra worker (tre soli campioni per
configurazione) è dello stesso ordine di grandezza dell'effetto misurato.
Questo guadagno non è gratuito: il traffico cresce nella stessa proporzione
di $k$ (Sez.~\ref{sec:results}C), per cui la scelta di $k$ va sempre motivata
rispetto al budget di banda disponibile, non solo rispetto all'accuratezza
target.

\subsection{Comunicazione non è il collo di bottiglia}
Un risultato rilevante ai fini dell'efficienza di comunicazione è che, anche
su rete reale tra istanze EC2 fisicamente separate, la Fase C resta
un'operazione dell'ordine delle centinaia di millisecondi, mentre il
rallentamento di oltre un ordine di grandezza osservato su AWS rispetto al
locale è concentrato nella Fase A (Sez.~\ref{sec:results}D), compatibile con
un limite di CPU burstable delle istanze \texttt{t3.small} piuttosto che con
un costo di rete. Questo supporta la scelta progettuale di comunicare
snapshot completi del modello a bassa frequenza ($H$ passi locali tra un
invio e l'altro) invece di aggiornamenti più frequenti e piccoli: nel regime
qui testato il costo di calcolo locale domina comunque quello di
comunicazione, per cui ridurre ulteriormente la frequenza di gossip non
avrebbe portato benefici proporzionali sul tempo totale.

\subsection{Nota metodologica sulla metrica di convergenza funzionale}
Lo strumento di analisi usato calcola lo spread di accuratezza sul global
test set in due varianti: al round finale di ciascun worker (quella usata
nella documentazione tecnica interna del progetto, ma esplicitamente
etichettata come riferimento rumoroso) e al round del checkpoint migliore di
ciascun worker (la metrica dichiarata primaria). Le due varianti possono
discordare sensibilmente per lo stesso run: per due delle configurazioni a
$N=3$, la variante \emph{riferimento} suggeriva uno spread contenuto
(4,5--4,8\%), mentre la metrica primaria mostra per le stesse run uno spread
del 19,5--21,1\%, in linea con tutte le altre 14 configurazioni. In questo
lavoro si è scelto di riportare sempre e solo la metrica primaria
(Tabella~\ref{tab:master}), proprio per evitare di presentare come
rappresentativo un valore isolato e più favorevole ottenuto con la metrica
esplicitamente definita come rumorosa.

\subsection{Limitazioni}
\begin{itemize}
  \item \textbf{Convergenza funzionale incompleta.} In nessuna delle 16
        configurazioni lo spread di accuratezza sul global test set scende
        sotto il 14\%: il protocollo gossip riduce ma non elimina la
        divergenza tra i modelli dei worker nel budget di round testato. I
        dati grezzi mostrano inoltre oscillazioni marcate round-su-round per
        singoli worker (in un caso, l'accuratezza global-test di un worker
        scende di 28 punti percentuali tra il proprio checkpoint migliore e
        il round finale), sintomo di instabilità subito dopo un merge FedAvg
        con vicini divergenti.
  \item \textbf{Singola replica per configurazione.} Nessun run è stato
        ripetuto con seed diverso: la deviazione standard riportata è
        \emph{tra worker} dello stesso run, non \emph{tra run} della stessa
        configurazione, e non permette di separare con certezza l'effetto di
        un iperparametro dalla variabilità intrinseca del processo.
  \item \textbf{Confondimento in una configurazione AWS.} Il run
        \texttt{n8\_fn} su AWS usa \texttt{early\_stopping\_patience=7}
        anziché 10 come tutti gli altri 15 run, per un errore di
        configurazione: il confronto diretto di questa configurazione con le
        altre a $N=8$ va quindi interpretato con cautela.
  \item \textbf{Peer effettivamente contattati sotto il target.} Il numero
        medio di vicini contattati da un worker può scendere sotto il valore
        di $k$ configurato quando altri worker terminano in anticipo
        (early stopping locale) e si deregistrano, riducendo il pool di peer
        disponibili — un effetto emergente dell'asincronia tra worker, non
        un difetto del meccanismo di selezione.
  \item \textbf{Assenza di heartbeat.} Il sistema rileva un peer caduto solo
        al primo fallimento di invio verso di esso; un crash non pulito
        (SIGKILL, OOM) lascia una entry non valida nel discovery server fino
        a quel momento, senza un meccanismo di TTL attivo. Questa è una
        scelta deliberata, non una svista: un heartbeat periodico verso il
        discovery server introdurrebbe traffico di coordinamento che cresce
        con $N$ e con la frequenza del battito, in diretto contrasto con
        l'obiettivo di traffico ridotto che guida le altre scelte
        protocollari di questo lavoro (cache dei peer, k-push a singolo hop);
        il costo è stato giudicato non giustificato per un caso d'uso — il
        crash non pulito — fuori dal perimetro sperimentale di questo lavoro.
  \item \textbf{Statistiche di Batch Normalization non aggregate.} Solo i
        parametri appresi (scala e bias) sono mediati da FedAvg; le
        statistiche correnti di media e varianza restano locali a ciascun
        worker, una scelta nota in letteratura per introdurre un
        disallineamento tra i client dopo l'aggregazione~\cite{li2021fedbn}.
\end{itemize}

\subsection{Confronto con la letteratura}
FedAvg centralizzato su FEMNIST è tipicamente riportato in letteratura
attorno al 77--80\% di accuratezza di test~\cite{li2020fedprox}. Le
configurazioni qui testate raggiungono un'accuratezza di validazione più alta
(87--91\%) ma un'accuratezza global-test più bassa (75--85\%): il confronto
diretto non è però del tutto omogeneo, poiché la validazione qui misurata
avviene su campioni tenuti da parte dagli stessi scrittori visti in training,
mentre il global test set contiene scrittori mai visti da nessun worker — un
compito di generalizzazione più severo, e più vicino allo scenario reale in
cui il modello deve funzionare su utenti non rappresentati nell'addestramento.
